% Options for packages loaded elsewhere
\PassOptionsToPackage{unicode}{hyperref}
\PassOptionsToPackage{hyphens}{url}
%
\documentclass[
]{article}
\usepackage{amsmath,amssymb}
\usepackage{iftex}
\ifPDFTeX
  \usepackage[T1]{fontenc}
  \usepackage[utf8]{inputenc}
  \usepackage{textcomp} % provide euro and other symbols
\else % if luatex or xetex
  \usepackage{unicode-math} % this also loads fontspec
  \defaultfontfeatures{Scale=MatchLowercase}
  \defaultfontfeatures[\rmfamily]{Ligatures=TeX,Scale=1}
\fi
\usepackage{lmodern}
\ifPDFTeX\else
  % xetex/luatex font selection
\fi
% Use upquote if available, for straight quotes in verbatim environments
\IfFileExists{upquote.sty}{\usepackage{upquote}}{}
\IfFileExists{microtype.sty}{% use microtype if available
  \usepackage[]{microtype}
  \UseMicrotypeSet[protrusion]{basicmath} % disable protrusion for tt fonts
}{}
\makeatletter
\@ifundefined{KOMAClassName}{% if non-KOMA class
  \IfFileExists{parskip.sty}{%
    \usepackage{parskip}
  }{% else
    \setlength{\parindent}{0pt}
    \setlength{\parskip}{6pt plus 2pt minus 1pt}}
}{% if KOMA class
  \KOMAoptions{parskip=half}}
\makeatother
\usepackage{xcolor}
\usepackage[margin=1in]{geometry}
\usepackage{graphicx}
\makeatletter
\def\maxwidth{\ifdim\Gin@nat@width>\linewidth\linewidth\else\Gin@nat@width\fi}
\def\maxheight{\ifdim\Gin@nat@height>\textheight\textheight\else\Gin@nat@height\fi}
\makeatother
% Scale images if necessary, so that they will not overflow the page
% margins by default, and it is still possible to overwrite the defaults
% using explicit options in \includegraphics[width, height, ...]{}
\setkeys{Gin}{width=\maxwidth,height=\maxheight,keepaspectratio}
% Set default figure placement to htbp
\makeatletter
\def\fps@figure{htbp}
\makeatother
\setlength{\emergencystretch}{3em} % prevent overfull lines
\providecommand{\tightlist}{%
  \setlength{\itemsep}{0pt}\setlength{\parskip}{0pt}}
\setcounter{secnumdepth}{-\maxdimen} % remove section numbering
\ifLuaTeX
  \usepackage{selnolig}  % disable illegal ligatures
\fi
\IfFileExists{bookmark.sty}{\usepackage{bookmark}}{\usepackage{hyperref}}
\IfFileExists{xurl.sty}{\usepackage{xurl}}{} % add URL line breaks if available
\urlstyle{same}
\hypersetup{
  pdftitle={Social Signaling and Prosocial Behavior},
  pdfauthor={Anne Karing (UChicago) and Karim Naguib ()},
  hidelinks,
  pdfcreator={LaTeX via pandoc}}

\title{Social Signaling and Prosocial Behavior}
\usepackage{etoolbox}
\makeatletter
\providecommand{\subtitle}[1]{% add subtitle to \maketitle
  \apptocmd{\@title}{\par {\large #1 \par}}{}{}
}
\makeatother
\subtitle{Experimental Evidence in Community Deworming in Kenya}
\author{Anne Karing (UChicago) and Karim Naguib ()}
\date{12/10/2020}

\begin{document}
\maketitle

\hypertarget{component-checks}{%
\subsection{Component checks}\label{component-checks}}

\hypertarget{fit-loading}{%
\subsection{Fit Loading}\label{fit-loading}}

\hypertarget{delta-expectation}{%
\paragraph{Delta Expectation}\label{delta-expectation}}

\begin{figure}
\includegraphics[width=\maxwidth]{takeup-STRUCTURAL_LINEAR_U_SHOCKS_PHAT_MU_REP89-fig/delta-with-u-example-1.pdf}
\caption{Net reputational return, $\mu\Delta[w^*]$, as a function of cutoff type, $w^*$}
\textit{Note:} {\footnotesize This plot shows how expected reputational return, $\mu\Delta[w^*]$ varies as cutoff type, $w^*$, changes. As $\sigma_u$ increases, the reputational return falls as it becomes harder to infer someone's type based off actions.}
\end{figure}

\begin{verbatim}
## [1] 4
\end{verbatim}

\hypertarget{reported-social-perception}{%
\paragraph{Reported Social
Perception}\label{reported-social-perception}}

\begin{figure}
\includegraphics[width=\maxwidth]{takeup-STRUCTURAL_LINEAR_U_SHOCKS_PHAT_MU_REP89-fig/reported-social-perception-1.pdf}
\caption{Reported social perception of some observable activities.}
\textit{Note:} {\footnotesize This plot shows the praise and stigma associated with various activities including deworming.}
\end{figure}

\hypertarget{distance-randomisation}{%
\paragraph{Distance Randomisation}\label{distance-randomisation}}

\begin{figure}
\includegraphics[width=\maxwidth]{takeup-STRUCTURAL_LINEAR_U_SHOCKS_PHAT_MU_REP89-fig/dist-plot-1.pdf}
\caption{Distance to Treatment Location (meters)}
\textit{Note:} {\footnotesize This graph shows a density plot of distance to the assigned point of treatment from each cluster centroid and household location. Clusters assigned to the close condition have greater density <1.25km whilst those assigned to the far condition place greater density >1.25km. Density estimates use all 9,805 surveyed individuals in the takeup sample.}
\end{figure}

\begin{figure}
\includegraphics[width=\maxwidth]{takeup-STRUCTURAL_LINEAR_U_SHOCKS_PHAT_MU_REP89-fig/dist-hist-1.pdf}
\caption{Distance to Treatment Location (meters)}
\textit{Note:} {\footnotesize This graph shows a histogram plot of distance to the assigned point of treatment from each cluster centroid and household location. Clusters assigned to the close condition have greater density <1.25km whilst those assigned to the far condition place greater density >1.25km. Density estimates use all 9,805 surveyed individuals in the takeup sample.}
\end{figure}

\begin{figure}
\includegraphics[width=\maxwidth]{takeup-STRUCTURAL_LINEAR_U_SHOCKS_PHAT_MU_REP89-fig/pot-distance-overlap-1.pdf}
\caption{Distribution of Distance to Closest Possible Point of Treatment, by Distance Condition}
\textit{Note:} {\footnotesize This plot shows there's significant overlap between the distribution of closest points of treatment across distance condition - i.e. any village assigned to a far PoT could have been assigned to a close PoT and vice versa.}
\end{figure}

\hypertarget{belief-facet-plots}{%
\paragraph{Belief Facet Plots}\label{belief-facet-plots}}

\begin{figure}
\includegraphics[width=\maxwidth]{takeup-STRUCTURAL_LINEAR_U_SHOCKS_PHAT_MU_REP89-fig/fob-beliefs-line-plot-post-1.pdf}
\caption{First Order Belief Levels and Average Treatment Effects}
\textit{Note:} {\footnotesize This plot shows first estimates using the first order beliefs  binomial model. Total sample size: 999 observations at the individual level.}
\end{figure}

\begin{figure}
\includegraphics[width=\maxwidth]{takeup-STRUCTURAL_LINEAR_U_SHOCKS_PHAT_MU_REP89-fig/sob-beliefs-line-plot-post-1.pdf}
\caption{Second Order Belief Levels and Average Treatment Effects}
\textit{Note:} {\footnotesize This plot shows first estimates using the second order beliefs  binomial model. Total sample size: 999 observations at the individual level.}
\end{figure}

\begin{verbatim}
## [1] 2
\end{verbatim}

\hypertarget{beliefs-plots}{%
\paragraph{Beliefs Plots}\label{beliefs-plots}}

\begin{figure}
\includegraphics[width=\maxwidth]{takeup-STRUCTURAL_LINEAR_U_SHOCKS_PHAT_MU_REP89-fig/belief-1ord-ate-1.pdf}
\caption{First Order Belief Average Treatment Effects}
\textit{Note:} {\footnotesize This plot shows first estimates using the first order beliefs  binomial model. Total sample size: 999 observations at the individual level.}
\end{figure}

\begin{figure}
\includegraphics[width=\maxwidth]{takeup-STRUCTURAL_LINEAR_U_SHOCKS_PHAT_MU_REP89-fig/belief-2ord-ate-1.pdf}
\caption{Second Order Belief Average Treatment Effects}
\textit{Note:} {\footnotesize This plot shows first estimates using the second order beliefs  binomial model. Total sample size: 999 observations at the individual level.}
\end{figure}

\begin{figure}
\includegraphics[width=\maxwidth]{takeup-STRUCTURAL_LINEAR_U_SHOCKS_PHAT_MU_REP89-fig/belief-1ord-prop-1.pdf}
\caption{First Order Belief Levels}
\textit{Note:} {\footnotesize This plot shows first estimates using the first order beliefs  binomial model. Total sample size: 999 observations at the individual level.}
\end{figure}

\begin{figure}
\includegraphics[width=\maxwidth]{takeup-STRUCTURAL_LINEAR_U_SHOCKS_PHAT_MU_REP89-fig/belief-2ord-prop-1.pdf}
\caption{Second Order Belief Levels}
\textit{Note:} {\footnotesize This plot shows first estimates using the second order beliefs  binomial model. Total sample size: 999 observations at the individual level.}
\end{figure}

\hypertarget{belief-ates-combined}{%
\paragraph{Belief ATEs combined}\label{belief-ates-combined}}

\begin{figure}
\includegraphics[width=\maxwidth]{takeup-STRUCTURAL_LINEAR_U_SHOCKS_PHAT_MU_REP89-fig/belief-1ord-prop-combined-1.pdf}
\caption{First Order Belief Levels - Combined}
\textit{Note:} {\footnotesize This plot shows estimates using the first order beliefs  binomial model. Estimates are combined across the close and far condition.  Total sample size: 999 observations at the individual level.}
\end{figure}

\begin{figure}
\includegraphics[width=\maxwidth]{takeup-STRUCTURAL_LINEAR_U_SHOCKS_PHAT_MU_REP89-fig/belief-2ord-prop-combined-1.pdf}
\caption{Second Order Belief Levels - Combined}
\textit{Note:} {\footnotesize This plot shows estimates using the second order beliefs  binomial model. Estimates are combined across the close and far condition.  Total sample size: 999 observations at the individual level.}
\end{figure}

\begin{figure}
\includegraphics[width=\maxwidth]{takeup-STRUCTURAL_LINEAR_U_SHOCKS_PHAT_MU_REP89-fig/belief-1ord-ate-combined-1.pdf}
\caption{First Order Belief Average Treatment Effects - Combined}
\textit{Note:} {\footnotesize This plot shows estimates using the first order beliefs  binomial model. Estimates are combined across the close and far condition.  Total sample size: 999 observations at the individual level.}
\end{figure}

\begin{figure}
\includegraphics[width=\maxwidth]{takeup-STRUCTURAL_LINEAR_U_SHOCKS_PHAT_MU_REP89-fig/belief-2ord-ate-combined-1.pdf}
\caption{Second Order Belief Average Treatment Effects - Combined}
\textit{Note:} {\footnotesize This plot shows estimates using the second order beliefs  binomial model. Estimates are combined across the close and far condition.  Total sample size: 999 observations at the individual level.}
\end{figure}

\hypertarget{disagg-beliefs}{%
\paragraph{Disagg Beliefs}\label{disagg-beliefs}}

\begin{verbatim}
## [1] 2
\end{verbatim}

\begin{figure}
\includegraphics[width=\maxwidth]{takeup-STRUCTURAL_LINEAR_U_SHOCKS_PHAT_MU_REP89-fig/disagg-fob-close-1.pdf}
\caption{Disaggregated First Order Belief Levels - Close}
\textit{Note:} {\footnotesize This plot shows estimates using the first-order beliefs  simple means.  Total sample size: 999 observations at the individual level.}
\end{figure}

\begin{verbatim}
## [1] 2
\end{verbatim}

\begin{figure}
\includegraphics[width=\maxwidth]{takeup-STRUCTURAL_LINEAR_U_SHOCKS_PHAT_MU_REP89-fig/disagg-fob-far-1.pdf}
\caption{Disaggregated First Order Belief Levels - Close}
\textit{Note:} {\footnotesize This plot shows estimates using the first-order beliefs  simple means.  Total sample size: 999 observations at the individual level.}
\end{figure}

\begin{verbatim}
## [1] 2
\end{verbatim}

\begin{figure}
\includegraphics[width=\maxwidth]{takeup-STRUCTURAL_LINEAR_U_SHOCKS_PHAT_MU_REP89-fig/disagg-fob-far-single-ci-1.pdf}
\caption{Disaggregated First Order Belief Levels - Close}
\textit{Note:} {\footnotesize This plot shows estimates using the first-order beliefs  simple means.  Total sample size: 999 observations at the individual level.}
\end{figure}

\begin{verbatim}
## [1] 2
\end{verbatim}

\begin{figure}
\includegraphics[width=\maxwidth]{takeup-STRUCTURAL_LINEAR_U_SHOCKS_PHAT_MU_REP89-fig/disagg-fob-both-1.pdf}
\caption{Disaggregated First Order Belief Levels}
\textit{Note:} {\footnotesize This plot shows estimates using the first-order beliefs  simple means.  Total sample size: 999 observations at the individual level.}
\end{figure}

\begin{verbatim}
## [1] 4
\end{verbatim}

\begin{figure}
\includegraphics[width=\maxwidth]{takeup-STRUCTURAL_LINEAR_U_SHOCKS_PHAT_MU_REP89-fig/disagg-sob-both-1.pdf}
\caption{Disaggregated Second Order Belief Levels}
\textit{Note:} {\footnotesize This plot shows estimates using the first-order beliefs  simple means.  Total sample size: 999 observations at the individual level.}
\end{figure}

\begin{verbatim}
## # A tibble: 4 x 7
##   assigned.treatment dist.pot.group dewormed gift_choice pref.~1 arm.s~2 assig~3
##   <fct>              <fct>          <lgl>    <fct>         <dbl>   <int> <chr>  
## 1 Calendar Arm       close          FALSE    Calendar      0.514     140 Calend~
## 2 Calendar Arm       far            FALSE    Calendar      0.629     159 Calend~
## 3 Bracelet Arm       close          FALSE    Bracelet      0.109     147 Bracel~
## 4 Bracelet Arm       far            FALSE    Bracelet      0.154     156 Bracel~
## # ... with abbreviated variable names 1: pref.prop[,1], 2: arm.size,
## #   3: assigned_treatment
\end{verbatim}

\begin{figure}
\includegraphics[width=\maxwidth]{takeup-STRUCTURAL_LINEAR_U_SHOCKS_PHAT_MU_REP89-fig/pref-gift-plot-1.pdf}
\caption{Preferences for gift by treatment arm}
\textit{Note:} {\footnotesize This plot shows the preference for either the bracelet or calendar gift by each treatment arm/distance condition.}
\end{figure}

\hypertarget{incentive-ates}{%
\paragraph{Incentive ATEs}\label{incentive-ates}}

\begin{figure}
\includegraphics[width=\maxwidth]{takeup-STRUCTURAL_LINEAR_U_SHOCKS_PHAT_MU_REP89-fig/reduced-incentive-level-plot-1.pdf}
\caption{Incentive Takeup Levels, Reduced Form Model}
\textit{Note:} {\footnotesize This plot shows the level of takeup in each treatment condition by distance, estimated using a Bayesian probit ('reduced form') model. The prior predictive distribution is overlaid at the top of the plot, which shows the predicted treatment effects before conditioning on the data. The prior is centered at 50\% takeup and the 90\% credible interval covers almost the entire unit interval. The plot shows that deworming takeup is higher in the close condition and that the bracelet has the highest level amongst the offered treatments. Line range: 90\% credible interval. Outer box: 80\% credible interval. Inner box: 50\% credible interval. Point: mean. 9,805 observations from the individual takeup sample used for estimates.}
\end{figure}

\begin{figure}
\includegraphics[width=\maxwidth]{takeup-STRUCTURAL_LINEAR_U_SHOCKS_PHAT_MU_REP89-fig/reduced-incentive-ate-plot-1.pdf}
\caption{Incentive Average Treatment Effect, Reduced Form Model}
\textit{Note:} {\footnotesize Line range: 90\% credible interval. Outer box: 80\% credible interval. Inner box: 50\% credible interval. Thick vertical line: median. Point: mean.}
\end{figure}

\begin{verbatim}
## [1] 2
\end{verbatim}

\begin{figure}
\includegraphics[width=\maxwidth]{takeup-STRUCTURAL_LINEAR_U_SHOCKS_PHAT_MU_REP89-fig/structural-signalling-ate-plot-1.pdf}
\caption{Signalling Average Treatment Effect, Structural Model}
\textit{Note:} {\footnotesize This plot shows the average signalling treatment effect estimated using the structural model. The prior predictive distribution is overlaid at the top of the plot, which shows the predicted signalling treatment effects before conditioning on the data. The prior is centered at a treatment effect of approximately 0. The plot shows that the signalling effect is greatest for the ink and bracelet treatments and largest in the far condition. Line range: 90\% credible interval. Outer box: 80\% credible interval. Inner box: 50\% credible interval. Thick vertical line: median. Point: mean. Signalling effect calculated using structural model incorporating willingness-to-pay, beliefs, and takeup data.}
\end{figure}

\begin{verbatim}
## [1] 2
\end{verbatim}

\begin{figure}
\includegraphics[width=\maxwidth]{takeup-STRUCTURAL_LINEAR_U_SHOCKS_PHAT_MU_REP89-fig/structural-private-ate-plot-1.pdf}
\caption{Private Incentive Average Treatment Effect, Structural Model}
\textit{Note:} {\footnotesize This plot shows the average private incentive treatment effect estimated using the structural model. The prior predictive distribution is overlaid at the top of the plot, which shows the predicted priavte incentive treatment effects before conditioning on the data. The prior is centered at a treatment effect of approximately 0. The plot shows that the private incentive effect is greatest for the calendar and bracelet treatments. Line range: 90\% credible interval. Outer box: 80\% credible interval. Inner box: 50\% credible interval. Thick vertical line: median. Point: mean.}
\end{figure}

\begin{figure}
\includegraphics[width=\maxwidth]{takeup-STRUCTURAL_LINEAR_U_SHOCKS_PHAT_MU_REP89-fig/structural-private-ate-plot-2.pdf}
\caption{Private Incentive Average Treatment Effect, Structural Model}
\textit{Note:} {\footnotesize This plot shows the average private incentive treatment effect estimated using the structural model. The prior predictive distribution is overlaid at the top of the plot, which shows the predicted priavte incentive treatment effects before conditioning on the data. The prior is centered at a treatment effect of approximately 0. The plot shows that the private incentive effect is greatest for the calendar and bracelet treatments. Line range: 90\% credible interval. Outer box: 80\% credible interval. Inner box: 50\% credible interval. Thick vertical line: median. Point: mean.}
\end{figure}

\begin{verbatim}
## [1] 2
\end{verbatim}

\begin{figure}
\includegraphics[width=\maxwidth]{takeup-STRUCTURAL_LINEAR_U_SHOCKS_PHAT_MU_REP89-fig/structural-incentive-ate-plot-1.pdf}
\caption{Incentive Average Treatment Effect, Structural Model}
\textit{Note:} {\footnotesize This plot shows the average treatment effect estimated using the structural model. The prior predictive distribution is overlaid at the top of the plot, which shows the predicted treatment effects before conditioning on the data. The prior is centered at a treatment effect of 0 and the 90\% credible interval covers the interval $[-0.125,0.125]$. The plot shows that the treatment effect is greatest for the bracelet treatment and estimated treatment effects are larger for the bracelet at close than far. Line range: 90\% credible interval. Outer box: 80\% credible interval. Inner box: 50\% credible interval. Thick vertical line: median. Point: mean. Treatment effect calculated using structural model incorporating willingness-to-pay, beliefs, and takeup data.}
\end{figure}

\begin{figure}
\includegraphics[width=\maxwidth]{takeup-STRUCTURAL_LINEAR_U_SHOCKS_PHAT_MU_REP89-fig/wtp-post-1.pdf}
\caption{Willingness-To-Pay Estimates}
\textit{Note:} {\footnotesize Calendar valued at 47.2KSh more than Bracelet with a 90\% credibility interval of $[41.2, 52.8]$. Line range: 90\% credible interval. Outer box: 80\% credible interval. Inner box: 50\% credible interval. Thick vertical line: median. Point: mean.  Preferences estimated using 998 individuals with a Bayesian multinomial discrete choice model using a Probit likelihood.}
\end{figure}

\begin{verbatim}
## # A tibble: 1 x 3
##   valuation_diff_estimate conf.low conf.high
##                     <dbl>    <dbl>     <dbl>
## 1                    46.8     41.7      53.1
\end{verbatim}

\begin{verbatim}
## [1] 2
\end{verbatim}

\begin{figure}
\includegraphics[width=\maxwidth]{takeup-STRUCTURAL_LINEAR_U_SHOCKS_PHAT_MU_REP89-fig/wtp-post-joint-1.pdf}
\caption{Willingness-To-Pay Estimates}
\textit{Note:} {\footnotesize Calendar valued at 47.2KSh more than Bracelet with a 90\% credibility interval of $[41.2, 52.8]$. Line range: 90\% credible interval. Outer box: 80\% credible interval. Inner box: 50\% credible interval. Thick vertical line: median. Point: mean.  Preferences estimated using 998 individuals with a Bayesian multinomial discrete choice model using a Probit likelihood.}
\end{figure}

\hypertarget{distance-prob-plots}{%
\paragraph{Distance Prob Plots}\label{distance-prob-plots}}

\begin{figure}
\includegraphics[width=\maxwidth]{takeup-STRUCTURAL_LINEAR_U_SHOCKS_PHAT_MU_REP89-fig/take-up-prob-with-obs-takeup-and-linerange-1.pdf}
\caption{Structural Model Estimates of Take-up Probability by Distance}
\textit{Note:} {\footnotesize This plot shows estimated take-up probability across bracelet and control over distance. Take-up probability is calculated using the structural model's estimates of $w^*$ in each cluster and taking the complement of the Gaussian CDF evalutaed at $w^*$. The sample consists of TODO. Outer ribbon: 90\% credibility interval, inner ribbon: 50\% credibility interval.}
\end{figure}

\hypertarget{rep-return-by-distance}{%
\subsection{Rep Return By Distance}\label{rep-return-by-distance}}

\begin{figure}
\includegraphics[width=\maxwidth]{takeup-STRUCTURAL_LINEAR_U_SHOCKS_PHAT_MU_REP89-fig/rep-return-dist-value-1.pdf}
\caption{Value of Reputational Return Across Bracelet and Control}
\textit{Note:} {\footnotesize Value priced in terms of walking distance. Estimates generated using decomposition from structural model.}
\end{figure}

\hypertarget{social-multiplier}{%
\subsection{Social Multiplier}\label{social-multiplier}}

\begin{figure}
\includegraphics[width=\maxwidth]{takeup-STRUCTURAL_LINEAR_U_SHOCKS_PHAT_MU_REP89-fig/rate-of-change-1.pdf}
\caption{Rate of Change Estimates}
\textit{Note:} {\footnotesize This plot shows the rate of change estimated using the structural model. The plot shows how sensitive take-up is to changes in cost - as distance increases individuals become less and less sensitive to the increased cost due to greater associated reputational returns. At 2.5km individuals in the bracelet condition are almost indifferent to an increase in costs (social multiplier close to 0). Line: median. Outer ribbon: 90\% credible interval. Inner ribbon: 50\% credible interval.}
\end{figure}

\hypertarget{sms}{%
\subsection{SMS}\label{sms}}

\begin{figure}
\includegraphics[width=\maxwidth]{/home/ed/projects/takeup/temp-data/sms-noreminder/TE-by-sms-incentive.pdf}
\caption{Average Treatment Effect, with and without SMS}
\textit{Note:} {\footnotesize }
\end{figure}

\begin{figure}
\includegraphics[width=\maxwidth]{/home/ed/projects/takeup/temp-data/sms-noreminder/TE-diff-by-sms-incentive.pdf}
\caption{Difference in Average Treatment Effect, Across SMS Treatment and Control}
\textit{Note:} {\footnotesize }
\end{figure}

\end{document}
